% !PRINTED 10
\weberpage{10}

%<<CONT
and among them one can select two in each case, for which the interval $z = 0$ to $z = 1$ corresponds to a given one of the intervals $(x_1x_2)$, $(x_2x_3)$, $(x_3x_4)$, $(x_4x_1)$; to increasing $z$ there correspond, in one of these transformations, the increasing $x$, and in the other the decreasing $x$ (with the passage from $+\infty$ to $-\infty$). In the first case the square roots in (9) have the same sign on both sides, in the second the opposite sign.

If one does not insist that $\varkappa^2$ be a positive proper fraction, then the points $x_1$, $x_2$, $x_3$, $x_4$ may be permuted in all ways, and one obtains 24 different linear transformations into the normal form, of which four at a time yield the same $\varkappa^2$; one obtains in all six different moduli, which are derived from one another according to § 2, (6).

One surveys the totality of these transformations most easily by assuming that the integral to be transformed already has the normal form:

\[
\dfrac{dx}{\sqrt{x(1-x)(1-\lambda^2x)}};
\]

one has then to replace, in the formulae (6) to (9), the $x_1$, $x_2$, $x_3$, $x_4$ in all possible ways by $0$, $1$, $1:\lambda^2$, $\infty$. We take as an example the following assignments:

\[
\begin{array}{ccccc}
 & x_1, & x_2, & x_3, & x_4, \\\hline
1) & 0, & \infty, & \dfrac{1}{\lambda^2}, & 1, \\[10pt]
2) & 0, & \dfrac{1}{\lambda^2}, & 1, & \infty, \\[10pt]
3) & 1, & 0, & \infty, & \dfrac{1}{\lambda^2}.
\end{array}
\]

One then obtains from (6) and (8)

\[
\begin{array}{c@{\quad}ll}
1) & z = \dfrac{-x}{1-x}, & \varkappa^2 = 1-\lambda^2 = \lambda'^2, \\[10pt]
2) & z = \lambda^2x, & \varkappa^2 = \dfrac{1}{\lambda^2}, \\[10pt]
3) & z = \dfrac{1-x}{1-\lambda^2x}, & \varkappa^2 = \lambda^2
\end{array}
\]

and for the differential

\[
\dfrac{dz}{\sqrt{z(1-z)(1-\varkappa^2z)}}
\]
%CONT>>
